\subsection{Rigid Fields}

% \subsection{Plane Rigid Fields}

A rigid body motion \emph{within the cross-sectional plane}
\((y,z)\) is a 3-parameter family,

\[
\bm{u}^{\text{plane rigid}} 
= \begin{pmatrix} 0 \\ \delta_y \\ \delta_z \end{pmatrix} 
+ \delta_\vartheta \begin{pmatrix}
0 \\ -z \\ y 
\end{pmatrix}
= \bm{\delta}_v + \delta_{\vartheta}(\FrameBasis \times \bm{r})
\]
where
\(\bm{\delta}_v \in \FrameBasis^\perp\)
represents in-plane translation and \(\delta_\vartheta \in \mathbb{R}\)
represents rotation about the \(\FrameBasis\)-axis. 

\subsection{Ieșan's Postulated Forms}

\subsubsection{Problem I}

In Problem I, Ieșan postulates that \(\hat{\bm{u}}'\) is rigid on each
section:
\[
\hat{\bm{u}}'_{\text{I}} = a_\varepsilon \FrameBasis + (a_\vartheta \FrameBasis - \FrameBasis \times \bm{a}_{\Section}) \times \bm{r}
\tag{V2}
\]
where \(a_\varepsilon \in \mathbb{R}\) is the axial strain,
\(a_\vartheta \in \mathbb{R}\) is the twist rate, and
\(\bm{a}_{\Section} \in \mathbb{R}^2\) are the bending
curvatures.

This gives the displacement:
\[
\hat{\bm{u}} = \bm{c}_0 + \bm{c}_1 \times \bm{r} + \reflenx\,(a_\varepsilon \FrameBasis + (a_\vartheta \FrameBasis - \FrameBasis \times \bm{a}_{\Section}) \times \bm{r}) + \text{(warping terms)}
\tag{V1}
\]

where \(\bm{c}_0 \in \mathbb{R}^3\) and
\(\bm{c}_1 \in \mathbb{R}^3\) are integration constants representing
the rigid body motion at \(x=0\).


\[
\hat{\bm{u}}_{\text{I}} = \bm{c}_0 + \bm{c}_1 \times \bm{r} + x[\text{prescribed } \hat{\bm{u}}'_{\text{I}}] + \text{(warping)}
\tag{V3}
\]
%
This introduces 6 integration constants:
\(\bm{c}_0 \in \mathbb{R}^3\) and \(\bm{c}_1 \in \mathbb{R}^3\).


\subsubsection{Problem II}
The general solution includes:
\[
\hat{\bm{u}}_{\rm I\!I} = \bm{c}_0 + \bm{c}_1 \times \bm{r} + x(\bm{c}_2 + \bm{c}_3 \times \bm{r}) + \text{(particular solution)}
\tag{V1}
\]

However, when Ieșan says \(\tilde{w}_\alpha = 0\) ``modulo a plane rigid
displacement,'' he means the transverse components of the warping can be
written as:
\[\tilde{\mathbf{w}}_\perp = \mathbf{d}_\perp + d_\vartheta(\FrameBasis \times \bm{r})\]

where \(\mathbf{d}_\perp \in \FrameBasis^\perp\) and
\(d_\vartheta \in \mathbb{R}\) are independent of
\(\bm{r}\).


\subsubsection{General Solution}

every term in Ieșan's general solution falls into one of three categories:

1. **Rigid body motions** of the section (absorbed into $\bm{u} + \bm{\theta} \times \bm{r}$)
2. **Warping shapes** that are $x$-independent functions of $\bm{r}$ (the columns of $\bm{\Phi}$)
3. **Products** of powers of $x$ with the above

\subsection{Extras}
The complete Saint~Venant solution consists of a particular solution satisfying the field equations plus an arbitrary rigid body displacement. 
For the combined extension-bending-torsion-flexure problem, the general solution takes the form:

\begin{equation}
\hat{\bm{u}}(x,\bm{r}) = \hat{\bm{u}}_{\rm P}(x,\bm{r}) + \hat{\bm{u}}_{\rm R}(x,\bm{r})
\end{equation}

where the rigid body field contains six arbitrary parameters \(\bm{\delta}\) representing the freedom to choose the reference configuration (TODO: Only the reference configuration?). 
Following the integration of Ie\c{s}an's postulated forms, this rigid field can be written as:

\begin{equation}
\hat{\bm{u}}_{\rm R} = \bm{\delta}_u + \bm{\delta}_\theta \times \bm{r} + x(\bm{\delta}_v + \bm{\delta}_\omega \times \bm{r})
\end{equation}

Here $\bm{\delta}_u, \bm{\delta}_v \in \mathbb{R}^3$ represent translational freedoms and $\bm{\delta}_\theta, \bm{\delta}_\omega \in \mathbb{R}^3$ represent rotational freedoms. For the pure flexure problem with vanishing end resultants at $x=0$, the constraints $\bm{\delta}_v \cdot \FrameBasis = 0$ and $\bm{\delta}_\omega = \mathbf{0}$ apply, reducing the arbitrary parameters.



\begin{Quote}
For the rotational freedom $\bm{q}_\theta$, we make the decomposition:
%
\begin{equation*}
\bm{q}_\theta = \delta_\vartheta\FrameBasis - \FrameBasis \times \bm{q}_{\Section}
\end{equation*}
%
where $\delta_\vartheta$ represents twist about the beam axis and $\bm{q}_{\Section} \in \mathbb{R}^2$ represents transverse rotation. 
\end{Quote}
\begin{Quote}
Following the integration of Ie\c{s}an's postulated forms, this rigid field can be written as:
%
\begin{equation*}
\hat{\bm{u}}_{\rm R} = \bm{q}_u + \bm{q}_\theta \times \bm{r} + x(\bm{q}_v + \bm{q}_\omega \times \bm{r})
\end{equation*}
%
Here $\bm{q}_u, \bm{q}_v \in \mathbb{R}^3$ represent translational freedoms and $\bm{q}_\theta, \bm{q}_\omega \in \mathbb{R}^3$ represent rotational freedoms. 
\end{Quote}